\documentclass[11pt,twocolumn]{article}
\usepackage[utf8]{inputenc}
\usepackage[french]{babel}
\usepackage{geometry}
\usepackage{booktabs}
\usepackage{hyperref}
\usepackage{amsmath}

% Contrainte : Limite de 2 pages (excluant les références et contributions) 
\geometry{margin=2cm}

% Contrainte : Titre préfixé par le numéro de groupe 
\title{Groupe XX : Architecture Hybride pour la Détection de Fausses Nouvelles par Analyse Stylométrique et Sémantique}

% Contrainte : Uniquement les noms et IDs des membres 
\author{Prénom Nom \\ Prénom Nom}
\date{23 mars 2026}

\begin{document}

\maketitle

\section{Énoncé du problème}
La tâche consiste à classifier un texte comme vrai (1) ou faux (0). Cette détection est complexe car les fausses nouvelles imitent souvent le format journalistique tout en injectant des biais stylistiques pour manipuler l'opinion.
\textbf{Exemple :} 
\begin{itemize}
    \item \textit{Entrée :} "INCROYABLE : Ce fruit guérit le cancer en 24h ! Partagez avant que ce soit supprimé." 
    \item \textit{Sortie :} Fake News (Label 0). 
\end{itemize}

\section{Travaux connexes}
La littérature scientifique sur la détection de fausses nouvelles s'est historiquement concentrée sur deux piliers : l'analyse de contenu et la stylométrie.

\textbf{Analyse Stylométrique et POS-tagging :} Les travaux pionniers de \textit{Potthast et al. (2017)} sur le corpus "Hyperpartisan News" ont démontré que les textes trompeurs utilisent une structure syntaxique distincte. L'utilisation de l'étiquetage morphosyntaxique (POS-tagging) est cruciale : \textit{Pérez-Rosas et al. (2018)} ont observé que les fausses nouvelles présentent souvent un ratio plus élevé d'adverbes et d'adjectifs qualificatifs, visant à renforcer l'aspect émotionnel au détriment de la précision factuelle. Notre extraction de 25 caractéristiques s'inspire de ces découvertes pour capturer cette signature "subjective".

\textbf{Modèles de Langage et Hybridation :} Bien que les modèles basés sur l'architecture Transformer (comme RoBERTa) excellent dans la capture du contexte sémantique, ils sont parfois limités par le bruit stylistique des réseaux sociaux (hashtags, majuscules). L'approche hybride, combinant vecteurs sémantiques et métadonnées artisanales, est une technique émergente visant à compenser l'aspect "boîte noire" des réseaux de neurones profonds par des indicateurs linguistiques interprétables.

\section{Méthodologie}
Notre architecture en "Y" traite le texte en deux flux parallèles :
1. \textbf{Flux Sémantique :} RoBERTa extrait une probabilité de fausseté ($roberta\_proba$).
2. \textbf{Flux Stylistique :} Un module \textit{StyleExtractor} calcule 25 métriques (ponctuation agressive, ratio de majuscules, POS ratios, sentiment via TextBlob).

\textbf{Fusion et Modèles de Classification :} Pour fusionner ces métriques stylistiques avec la probabilité de RoBERTa, nous avons mis en compétition deux modèles via une recherche d'hyperparamètres (\textit{RandomizedSearchCV}):
1. \textbf{XGBoost :} Reconnu pour ses performances sur les données tabulaires.
2. \textbf{Random Forest :} Reconnu pour sa robustesse et sa capacité à réduire la variance, qui a finalement été sélectionné comme notre juge final.

\section{Expériences et Évaluation}
\subsection{Données}
Nous fusionnons le dataset \textbf{LIAR} (déclarations politiques) et \textbf{Disaster Tweets}. Les labels ont été binarisés et les entrées "half-true" exclues pour garantir la clarté de l'apprentissage.

\subsection{Base de référence (Baseline)}
Notre baseline est l'utilisation du modèle \textbf{RoBERTa seul} (via sa probabilité de sortie $roberta\_proba$, sans les 25 caractéristiques de style). Cela nous permet de mesurer précisément la valeur ajoutée de notre extracteur stylistique artisanal par rapport à une approche purement sémantique Deep Learning.

\section{Résultats préliminaires}
\subsection{Performance Quantitative}
Les tests sur nos 12 335 échantillons démontrent que la baseline sémantique (RoBERTa) accomplit déjà une excellente classification. L'ajout du style affine ces résultats : le Random Forest s'impose légèrement sur la classification stricte (Accuracy), tandis que XGBoost améliore drastiquement la calibration de la confiance (Log Loss).

\begin{table}[h]
\centering
\begin{small}
\begin{tabular}{lcccc}
\toprule
Modèle & Accuracy & F1-score & ROC-AUC & Log Loss \\
\midrule
RoBERTa seul & 92.20\% & 91.98\% & 98.47\% & 16.75\% \\
XGBoost & 92.18\% & 91.95\% & \textbf{98.50\%} & \textbf{15.11\%} \\
\textbf{Random Forest} & \textbf{92.24\%} & \textbf{92.03\%} & 98.44\% & 17.58\% \\
\bottomrule
\end{tabular}
\end{small}
\caption{Comparaison des performances des modèles hybrides avec la baseline.}
\end{table}

\textbf{Analyse détaillée du meilleur modèle (Random Forest) :}
Le modèle démontre un excellent équilibre entre le rappel et la précision pour les deux classes, avec une précision de 0.95 pour la classe 1 (Vraies Nouvelles) et un rappel de 0.95 pour la classe 0 (Fausses Nouvelles).

\subsection{Analyse de l'Importance (Feature Importance)}
L'analyse révèle que si la sémantique (\textit{roberta\_proba}) domine largement les décisions du Random Forest, les caractéristiques stylistiques jouent un rôle décisif pour affiner la prédiction et surpasser la baseline. La \textbf{densité de ponctuation agressive} et le \textbf{ratio d'adverbes} se détachent comme les prédicteurs stylistiques majeurs, confirmant que le ton sensationnaliste est une signature mathématiquement mesurable des fausses nouvelles.

% \section{Travaux futurs}
% Pour la suite du projet, nous prévoyons d'intégrer des indices de lisibilité (Gunning Fog) pour tester si la simplification syntaxique est corrélée aux fausses nouvelles, et d'optimiser le pipeline d'inférence pour la démonstration finale.

\bibliographystyle{acl_natbib}
\begin{thebibliography}{9}
\bibitem{potthast} Potthast, M., et al. (2017). Stylometric Inquiry into Hyperpartisan News Detection.
\bibitem{perez} Pérez-Rosas, V., et al. (2018). Automatic Detection of Fake News.
\end{thebibliography}

% \section{Contributions}
% \begin{itemize}
%     \item \textbf{Membre A :} Développement du \textit{StyleExtractor} (POS tagging, sentiment), binarisation des données et rédaction de la méthodologie.
%     \item \textbf{Membre B :} Entraînement de RoBERTa, optimisation des modèles XGBoost/RF et analyse de l'importance des variables.
% \end{itemize}

\end{document}